\section{Introduction}
Many NLP systems need a representation for a label. In some settings the label is represented by its name, by a natural-language definition, by a prompt, or by a small set of examples. This paper studies a simple but consequential case: a category is represented by aggregating embeddings of the instances assigned to it. We call these \emph{member-derived label representations}. Such representations are attractive because they require no new annotation beyond category membership. They are also risky: high semantic similarity among members may not imply that the resulting vector is useful for a downstream task.

We examine this risk through educational Knowledge Components (KCs). A KC is a human-defined label over educational questions, such as a math skill or concept. Each KC has member question texts, so a natural representation is the mean, median, medoid, or otherwise aggregated embedding of those questions. At the same time, KCs are not merely topical labels. In knowledge tracing (KT), they define the sequence of student--skill interactions used to predict future correctness. This dual role makes KCs a useful testbed for label representation: we can ask whether a representation recovers category membership, and separately whether it captures the functional structure needed by a predictive model.

The paper is motivated by a negative downstream result from Chapter 5 of the thesis. That chapter tested text-derived KC-response embeddings as initializations for Deep Knowledge Tracing (DKT; \citealp{piech2015dkt}). Across nine construction pipelines, semantic/category initialization did not reliably improve DKT over the standard learned interaction embedding. We reframe that result as an NLP label-representation problem. The relevant question is not whether text embeddings are ``good'' or ``bad'' in general. The question is whether a representation obtained from constituent members is valid for the function assigned to the label.

This short paper makes a focused point. Member-derived KC representations encode category information, but category validity is not the same as downstream functional validity. We show that prototypes strongly recover KC membership and partially align with heuristic label groupings. However, their alignment with functional KT graphs is weak, and Chapter 5 downstream initialization results remain mixed or negative. The contribution is therefore a scoped negative finding: member aggregation is a reasonable way to represent labels for category recovery, but it should not be assumed to produce a functional prior without task-specific evidence.

Our contributions are:
(1) we formalize KCs as human-defined text labels for studying member-derived label representation;
(2) we evaluate category recovery, category coherence, heuristic label alignment, functional graph alignment, difficulty prediction, and transition prediction;
(3) we consolidate Chapter 5 downstream DKT evidence as the main functional test; and
(4) we identify failure modes that explain why an embedding can represent category membership while failing as a KT initialization.
